ICAO defines runway incursions as any occurrence at an airport involving an aircraft,vehicle, person, or object on the ground that creates a collision hazard or results in loss of separation with an aircraft taking off, intending to take off, landing, or intending to land.". As the volume of air traffic rises, the risk of runway incursions at already congested airports increases. FAA reports that in 2023, there were 1,760 runway incursions in America, of which 1,070 incidents, or 60\% of the total, are attributed to pilot error. This metric only serves to highlight the need for a flight deck-based system to increase pilot's situational awareness in order to react to developing incidents. \newline

This study used historical runway incursion parameters gathered [incident reports] from FAA ASIAS, Skybrary, (insert international sources here), and international aviation regulatory bodies in order to recreate real world scenarios. Runway incursion events gathered included Category A and B incidents by the FAA's Runway Incursion Severity (RIS) model, as well as accidents. Parameters included type of involved aircraft, runway geometry, cause of the runway incursion, aircraft state, and severity rating assigned by the team in accordance to Operational Services and Environment Definition (OSED) aircraft states described within RTCA DO-323. The team then validated the severity rating by use of ICAO's RIS calculator before analysis.

%%As the RIS model is not adopted globally, non-US data entries were manually categorized by the team in accordance with the FAA RIS model (would be great if we could get our hands on FAA's RIS calculator). 

The team then employed the Surface Situational Awareness with Indications and Alerts (SURF-IA) detection logic on these scenarios and attempted to determine whether SURF IA would give a true alert, untrue alert, or whether it would miss the alert entirely in order to quanitfy the system's effectiveness on real world scenarios. \newline
[STUDY RESULTS HERE]